\section{Background And Related Work}

\subsection{Optical Design Workflows}

Optical lens design is commonly formulated as an iterative prescription refinement problem. A designer specifies an initial sequence of surfaces, materials, aperture stops, field conditions, and design targets; optical design software then evaluates the system through ray tracing and optimizes the available variables with respect to chosen merit functions \cite{Smith2008ModernOpticalEngineering,codev,zemax_opticstudio}. This software stack is highly effective once a viable initial prescription is available. However, early-stage exploration remains difficult because the merit landscape can be highly non-convex and sensitive to initialization \cite{Turnhout:09}. In many practical workflows, topology choices such as element count, surface ordering, material transitions, and stop placement are still guided by designer experience or by adapting existing designs.

\subsection{Generative Models For Engineering Design}

Deep generative models have become an increasingly important tool for exploring structured engineering design spaces \cite{regenwetter2022deepgenerativemodelsengineering}. Recent work has applied diffusion and latent generative models to CAD sequence generation, aircraft geometry, and ship hull design \cite{alam2025gencadimageconditionedcomputeraideddesign,blendednet_pp_Sung_2025,bagazinski2023shipgendiffusionmodelparametric}. These examples suggest that modern generative models can be useful when a design object has a meaningful representation and when generated samples can be evaluated against domain-specific constraints. They also motivate the central representation question in this work: whether lens prescriptions should be treated as generic tabular rows, structured surface-indexed tables, or ordered optical design objects.

\subsection{Learning-Based Optical Design}

Machine learning has previously been used in optical design for surrogate modeling, optimization acceleration, and starting-point generation. Surrogate models can accelerate repeated optical evaluations during optimization \cite{hegde2019accelerating}. Neural networks have also been used to generate starting points for specialized optical systems, such as freeform off-axis reflective triplets \cite{Yang:19}. More closely related to this work, recurrent neural networks have been used for lens-design starting-point generation \cite{Cote:21}. These studies show the potential of data-driven optical design, but existing lens-generation work has generally been limited by smaller datasets, restricted design families, fixed or manually specified topology assumptions, or simpler sequence models.

Lens prescriptions can be stored and generated as mixed numerical and categorical tables, which makes modern tabular generative models such as TabDDPM, STaSy, CoDi, and TabSyn a natural benchmark class for this work \cite{Kotelnikov2023TabDDPM,Kim2023STaSy,Lee2023CoDi,Zhang2024TabSyn}. However, a fixed raw prescription table does not explicitly tell the model that groups of columns belong to ordered surfaces along the optical axis, that adjacent surfaces are coupled, or that designs with fewer surfaces contain structural padding rather than ordinary missing data. This motivates a tuned tabular representation that keeps the data compatible with benchmark generative models while exposing optical scale, material-property, stop-location, and surface-count structure more directly. It also motivates a second, architecture-level representation in which prescriptions are treated as variable-length sequences of ordered surface-parameter tokens with explicit surface-position, parameter-type, and valid-surface information.
